\section{Problem Formulation}

We model an auditable clinical/research AI system as a set of typed authorities connected by governed transitions. The central design constraint is that semantically distinct objects must not silently inherit one another's authority:

\begin{equation}
\text{source} \neq \text{deterministic derivation} \neq \text{model inference} \neq \text{human-reviewed artifact} \neq \text{action proposal} \neq \text{external effect}.
\end{equation}

This distinction motivates four system properties.

\paragraph{Authority integrity.} A presentation layer, worker, model, projection, or integration must not become an alternate source of truth merely because it can produce or display data.

\paragraph{Provenance integrity.} A derived claim or artifact must retain sufficient typed linkage to determine what source identity and transformation it depends on, and integrity-sensitive bindings must fail when intentionally corrupted.

\paragraph{Epistemic honesty.} Unknown, unsupported, unavailable, denied, partial, conflicting, stale, corrupt, or unmeasured conditions must not be collapsed into a generic success state. The paper distinguishes vocabulary implemented directly in individual subsystems from a broader cross-system taxonomy proposed for analysis; it does not claim that every label is represented by one universal enum.

\paragraph{Effect integrity.} A model recommendation or internal action proposal must not be treated as evidence that an external action occurred. External effects require their own governed transition and evidence.

The paper evaluates these properties under synthetic inputs and an exact repository snapshot. Clinical correctness and patient outcomes are outside the current experimental scope.
